\documentclass[12pt]{article}
\usepackage[utf8]{inputenc}
\usepackage{amsmath,amssymb,amsthm}
\usepackage[margin=1in]{geometry}

\newtheorem{theorem}{Theorem}
\newtheorem{lemma}[theorem]{Lemma}

\title{Problem 155: Counting Capacitor Circuits}
\author{Project Euler Solution}
\date{}

\begin{document}
\maketitle

\section{Problem Statement}
Using up to 18 unit capacitors in series/parallel combinations, how many distinct capacitance values can be obtained? Define $D(n)$ as the set of values achievable with exactly $n$ capacitors. Find $C(18) = |D(1) \cup \cdots \cup D(18)|$.

\section{Mathematical Foundation}

\begin{theorem}[Recursive Construction]
$D(1) = \{1\}$ and for $n \ge 2$:
\[
D(n) = \bigcup_{k=1}^{\lfloor n/2 \rfloor} \left\{ a + b,\; \frac{ab}{a+b} \;\middle|\; a \in D(k),\; b \in D(n-k) \right\}
\]
\end{theorem}

\begin{proof}
Any circuit of $n \ge 2$ capacitors decomposes as a parallel or series combination of two sub-circuits using $k$ and $n-k$ capacitors. Parallel gives capacitance $a + b$; series gives $ab/(a+b)$. Since both operations are commutative, we restrict $k \le \lfloor n/2 \rfloor$. Completeness follows by structural induction on circuits.
\end{proof}

\begin{theorem}[Rationality]
Every element of $D(n)$ is a positive rational number.
\end{theorem}

\begin{proof}
By induction. $D(1) = \{1\} \subset \mathbb{Q}_{>0}$. If $a, b \in \mathbb{Q}_{>0}$, then $a + b \in \mathbb{Q}_{>0}$ and $ab/(a+b) \in \mathbb{Q}_{>0}$.
\end{proof}

\begin{lemma}[Series-Parallel Duality]
If $c \in D(n)$, then $1/c \in D(n)$.
\end{lemma}

\begin{proof}
By induction on circuit structure. Swapping every parallel connection with series and vice versa transforms capacitance $c$ to $1/c$, since parallel $a + b$ becomes series $ab/(a+b)$ of the reciprocal circuits (achieving $1/a$ and $1/b$), which gives $\frac{(1/a)(1/b)}{1/a + 1/b} = \frac{1}{a+b}$.
\end{proof}

\begin{lemma}[Canonical Representation]
Capacitances are stored as reduced fractions $(p, q)$ with $\gcd(p,q) = 1$ and $q > 0$. Equality of capacitances corresponds to equality of reduced fractions.
\end{lemma}

\begin{proof}
The reduced-fraction representation of a rational number is unique.
\end{proof}

\section{Algorithm}

\begin{verbatim}
D[1] = {1}
for n = 2 to 18:
    D[n] = {}
    for k = 1 to n/2:
        for (a, b) in D[k] x D[n-k]:
            D[n].insert(a + b)
            D[n].insert(a*b / (a+b))
return |D[1] union ... union D[18]|
\end{verbatim}

\section{Complexity Analysis}

\begin{itemize}
    \item \textbf{Time:} $\sum_{n=2}^{N} \sum_{k=1}^{\lfloor n/2 \rfloor} |D(k)| \cdot |D(n-k)|$ set insertions, each $O(1)$ amortized with hashing. Empirically $\sim 10^9$ operations for $N = 18$.
    \item \textbf{Space:} $O\!\left(\sum_{n=1}^{N} |D(n)|\right)$, on the order of millions of stored fractions.
\end{itemize}

\section{Answer}

\[
\boxed{3857447}
\]

\end{document}
